% Native bridge audit and the pre-selector finite-case repair.
% Requires amsmath, amssymb, amsthm and theorem/lemma environments.

\section{Native bootstrap: cases before a full selector}
\label{nb:section}

This section checks a dependency that matters for the new scalar basis.
Before predecessor is available, the partial selector
\[
 B(n,a,b)=J(n,a;u.b),\qquad B(0,a,b)=Sa,\quad B(Sn,a,b)=b
\]
is not the unrestricted selector $C(0,a,b)=a$, $C(Sn,a,b)=b$.
Accordingly, one must not silently invoke the later ordinary-connective
macros during the arithmetic bootstrap.  In particular, the theorem rule
for numerical cases is not by itself a rule for eliminating equations
$t=0$ and $t=Sz$ from a retained context.  The finite-case lemma below
supplies exactly the stronger form used by the bootstrap.

We use the native Eq/$Q$/$D$ presentation with hypothesis use, full
formula transport, $Q$ elimination, paired elimination, and the two
paired introductions.  In notation,
\[
 Q_yF=\text{the clause taking a proof of }F(0)\text{ to }F(t),
\]
\[
 \frac{\Gamma,F(t,0)\vdash F(t,Sz)}
      {\Gamma\vdash D_{x,y}F(x,y)},\qquad
 \frac{\Gamma,E(t,z)\vdash E(t,Sz)}
      {\Gamma\vdash D_{x,y}E(x,y)}.
\]
The second matrix is a literal equation.  The usual freshness conditions
are retained, and paired elimination opens $Q_yF(a,y)$ at a fresh $a$.
Computation equations may be read as leaves, or obtained from their
constant generated-assumption instances and hypothesis use.

\subsection{Resources available without a numerical selector}

Choose $v$ fresh for $t$ and put $r=J(S0,0;v.t)$.  Computation gives
$r=t$.  Transport this equation along itself with motive $w=t$ to obtain
$t=t$.  The familiar equality motives then give symmetry, transitivity
and numerical congruence; this uses neither predecessor nor successor
injectivity.

For any native formula template $F$, case introduction with an unused
witness coordinate gives $D_{x,y}F(y)$ from
$\Gamma\vdash F(Sz)$, ignoring the retained $F(0)$.  Paired elimination
opens $Q_yF(y)$; a supplied proof of $F(0)$ followed by $Q$ elimination
at $t$ gives
\[
 \frac{\Gamma\vdash F(0)\qquad\Gamma\vdash F(Sz)}
      {\Gamma\vdash F(t)}. \tag{Native NC}\label{nb:nc}
\]
The identical construction with progress introduction, using its local
$E(z)$, gives contextual literal-equation induction from a base and step.
These arguments use no arithmetic selector or ordinary logical connective.

For reference, $Q_yF$ and $D_{x,y}F$, with $x$ unused, are interderivable
already: use case introduction and $Q$ elimination in the forward
direction, and paired elimination followed by hypothesis use in the
reverse direction.  Thus the additional $Q$ constructor presents the
same internal clauses as the Eq/$D$ formulation.

Even equation explosion is available at this stage.  Under $0=1$,
congruence gives $B(0,0,t)=B(1,0,t)$; computation gives $1=t$.
Consequently every equation follows.  External induction on fixed native
formula syntax, using case introduction (and the unused-witness
representation for $Q$), gives explosion for every native formula.
This is an external proof transformation, not compound-formula induction
inside the arithmetic theory.

\subsection{Functional finite ranges do not assume branch equations}

Say that a particular term $b$ has the \emph{functional range rule}
$\operatorname{FR}_R(b)$, for a finite list of numerals $R$, when the
following metarule is derived for every native formula template $K$:
\[
 \bigl(\Gamma\vdash K(r)\text{ for each }r\in R\bigr)
 \quad\Longrightarrow\quad\Gamma\vdash K(b).
 \tag{FR}\label{nb:fr}
\]
There are no added equations $b=r$ in these premises.  The retained
context and every fixed parameter of $K$ remain unchanged.  In
particular, they may mention the original term $b$.

Use the fixed terms $Z,N,j,M$ defined in
Section~\ref{js:cycles}, with $Z(t)=B(t,0,0)$. The present proof
checks their logical prerequisites; their finite computation identities
are verified directly from those definitions.

For $Z(t)$, apply Native NC to the fresh template $w\mapsto K(Z(w))$.
The zero branch is $K(1)$ and every successor branch is $K(0)$, by
computation and transport.  Hence $\operatorname{FR}_{\{0,1\}}(Z(t))$
holds, and therefore so does the same rule for $N(t)$.
For $B(t,1,0)$ the identical construction has endpoints two and zero.
For
\[
 j(t)=B\bigl(t,0,B(Z(N(t)),1,0)\bigr)
\]
use NC on the outer control: the zero branch is one, and the successor
branch has the preceding two-value rule.  This proves
$\operatorname{FR}_{\{0,1,2\}}(j(t))$, and hence the same rule for
$M(t)$, since $M$ is an outer application of $j$.

The inner parity iterate $U(t)=J(t,0;u.Z(u))$ has a two-value range
rule as well.  At counter zero it is one; at a successor it is a value
of $Z$.  Apply NC on the counter and the already proved range rule of
$Z$.  No induction with a disjunctive invariant is involved.

These constructions justify the cycle equations before contextual
finite-case elimination is used.  For example, $Z(Z(U(t)))=U(t)$ follows
by applying its two-value rule to the equation template $Z(Z(w))=w$.
The equations
\[
 M(0)=0,\quad M(1)=1,\quad M(2)=2,\qquad M(Sw)=j(M(w))
\]
and $j(0)=1$, $j(1)=2$, $j(2)=0$ then follow from the fixed definitions.
Their functional range proofs are independent of the ramp invariant.

\begin{lemma}[Successor cancellation on the three finite values]
\label{nb:finite-cancellation}
Suppose $b$ and $c$ have functional range rules whose values are contained
in $\{0,1,2\}$.  In any retained context,
\[
                         Sb=Sc\quad\Longrightarrow\quad b=c
\]
is derivable, before global successor injectivity is known.
\end{lemma}
\begin{proof}
Apply the two functional range rules to the equation templates
$M(w)=w$ and $j(j(j(w)))=w$.  Their required numeral instances are fixed
computations.  Thus $M(b)=b$, $M(c)=c$, $j^3(b)=b$ and $j^3(c)=c$
are already theorems in the retained context.
Apply $M$ to $Sb=Sc$ and use $M(Sw)=j(M(w))$ to get $j(b)=j(c)$.
Apply $j$ twice and use the two cycle identities to obtain $b=c$.
Every step is ordinary numerical congruence or transport.  In particular,
no cancellation of an arbitrary successor was inserted.
\end{proof}

\begin{lemma}[Contextual finite-case elimination for equations]
\label{nb:guarded-cases}
Let $b$ have $\operatorname{FR}_R(b)$, where $R$ is contained in
$\{0,1,2\}$, and let $F$ be the literal equation $u=v$.
If, for each $r\in R$, there is a derivation
\[
                         \Gamma,b=r\vdash u=v,
\]
then $\Gamma\vdash u=v$ is derivable.  The context may contain arbitrary
native formulas, including an induction hypothesis that depends on $b$.
\end{lemma}
\begin{proof}
Define a native formula template, with both bound variables fresh,
\[
 K(x):=D_{a,y}\bigl(B(y,x,u)=B(y,b,v)\bigr),
\]
where $a$ is unused.  Fix $r\in R$.  To introduce $K(r)$ by the case
choice with witness zero, its local zero assumption reduces by
computation to $Sr=Sb$.  Lemma~\ref{nb:finite-cancellation}, with the
constant $r$ and the known finite-range term $b$, yields $b=r$.
Cut in the given branch derivation.  Its conclusion $u=v$ is exactly
the successor matrix, up to the computations of $B$.  Thus
$\Gamma\vdash K(r)$.

Apply the functional range rule of $b$ to these closed branch proofs.
It gives $\Gamma\vdash K(b)$, whose matrix is
$B(y,b,u)=B(y,b,v)$.  Paired elimination opens its $Q$ clause.
At zero its two sides are $Sb$, so reflexivity and computation supply
the base.  Specializing the clause at one gives $u=v$.  Close paired
elimination.  All temporary coordinates are fresh and all temporary
assumptions have been discharged.
\end{proof}

The derivation is a finite guard construction, not a primitive implication
connective.  It relies only on the partial selector, computations,
transport, hypotheses, case introduction, paired elimination and $Q$
elimination.  Neither predecessor nor the unrestricted selector occurs.

\subsection{Audit of the bootstrap and completion of the bridge}

The dependency order is now explicit.  First establish functional range
rules and the two finite cycles using only Native NC and equations.
Then establish finite successor cancellation and
Lemma~\ref{nb:guarded-cases}.  The finite splits in the ramp induction,
the affine scaling induction and the factor-map calculation are instances
of that lemma, with $b=N(n)$ or another displayed finite-range term and
with a literal-equation conclusion.  Retained induction hypotheses are
therefore legitimate.  The saturation identity $K(c,Su)=S K(c,u)$ is
proved directly by NC with no branch assumptions; it does not require
an unrestricted contextual split on arbitrary $c$.  No step in this order
uses the later ordinary logical macros.

After the corrected predecessor $P$ is constructed, the initializer
compiler gives the ordinary scalar $I$.  Consequently
\[
 C(n,a,b):=I(n,a;u.b)
\]
has the unrestricted selector endpoints.  The inherited gate, padding,
formula-selector and ordinary-connective constructions may now be applied
in their stated order.  Binder congruence itself uses only case
introduction, paired elimination, transport and cut; it is already valid
before this final stage.

For clarity, the ordinary interpretations of the native binders are
\[
 (Q_yB)^\circ=B^\circ(0)\to\forall y\,B^\circ(y),\qquad
 (D_{x,y}B)^\circ=
 \exists x\bigl(B^\circ(x,0)\to\forall y\,B^\circ(x,y)\bigr).
\]
Once the logical macros are available, the native return translations
are proved by the same explicit constructions:
\[
 Q_yB\dashv\vdash B(0)\Rightarrow_*\forall_*y\,B(y),
\]
\[
 D_{x,y}B\dashv\vdash
 \exists_*x\bigl(B(x,0)\Rightarrow_*\forall_*y\,B(x,y)\bigr).
\]
Forwards, an opened $Q$ clause under its base supplies the arbitrary
instance needed by universal introduction.  Backwards, a displayed
implication and universal supply the successor matrix for case
introduction; an unused witness gives the $Q$ case.  The existential
witness is retained in the $D$ case.  External structural induction on
finite formulas, together with the endpoint proofs for the ordinary
macros, yields both formula composites.  Falsity uses the explicit
ordinary separation bridge $(0=1)\to\bot$.

Finally, the inherited scalar completion and fixed-description PR
compiler are used only after the logical interface and initialized $I$
have been recovered.  Recurrence uniqueness gives the numerical
composites on the whole original domain.  The domain is unchanged,
interpreted equality is actual equality, and the numerical isomorphism
graph is $x=y$.  Thus the finite-case repair closes the pre-selector
logical dependency; it does not replace the separately supplied complete
arithmetic and logical interpretations by standard-model calculations.
